
\subsection{以外部資訊換取穩定：多母體、參考母體與修勻}\label{sec:ref}

第一類應對的共同想法是：既然困難出在資料的訊息量不足，
就從目標區域之外補進訊息。補進的方式有兩種，而兩種都以「外部資訊與目標
區域足夠相似」為前提。

其一是把母體合併。\citet{lilee2005} 合併死亡改善型態相近的多個母體，
以共同因子描述其共同趨勢；\citet{jarner2011} 以大型參考母體配適共同趨勢
而只讓小母體估計其偏離，此即 SAINT 模型；
\citet{ahcan2014} 把小母體的死亡資料與其他母體的資料混合，
以處理觀測序列中的劇烈隨機波動；
\citet{wang2018} 則建議聚合目標母體十至二十年的歷史資料。

其二是修勻。\citet{lee2003} 提出 partial SMR，以標準化死亡比把小區域的
年齡別死亡率朝參考母體收縮；\citet{yue2019} 把 partial SMR 與
\citet{whittaker1923} 修勻接在 Lee-Carter 之前或之後，比較兩種接法的表現。
修勻這一條值得展開，蓋因它與本文處理的是同一個病灶的兩種解法。
\citet{lee2003} 提出 partial SMR 的動機，與本文的動機幾乎相同：
該文在摘要中列出兩項困難，即小人口的年齡別率估計不穩定、
年齡曲線含過多隨機變異，以及「當分子含零計數時無法取對數」，
後者即零死亡格的困難。該文所選的解法是引入一個年齡曲線相近的標準母體，
以標準化死亡比為中心把觀測朝該母體收縮。

記 $d_x$ 為目標小區域 $x$ 歲組的死亡數觀測值、$P_x$ 為其曝露數、
$m^R_x$ 為參考母體同年齡組的中央死亡率，則期望死亡數為 $e_x=P_x m^R_x$，
標準化死亡比為 $\mathrm{SMR}=\sum_x d_x\big/\sum_x e_x$。
修勻後的死亡率定義為
\begin{equation}\label{eq:psmr}
v_x \;=\; m^R_x\,\exp\!\left(
\frac{d_x\hat h^2\log(d_x/e_x)+\bigl(1-d_x/\textstyle\sum_x d_x\bigr)\log \mathrm{SMR}}
     {d_x\hat h^2+\bigl(1-d_x/\textstyle\sum_x d_x\bigr)}\right),
\end{equation}
其中 $\hat h^2=\max\bigl\{[\sum_x(d_x-e_x\mathrm{SMR})^2-\sum_x d_x]
\big/[\mathrm{SMR}^2\sum_x e_x^2],\,0\bigr\}$ 度量兩個母體之間年齡別死亡率
比值的異質性。該文對 $\hat h^2$ 的構造說明值得一提：比值的異質性必然伴隨
超出 Poisson 所能解釋的變異，故可由觀測與期望死亡數的平方差減去 Poisson
本身的變異後估得，此即上式分子中減去 $\sum_x d_x$ 的來由。
式~\eqref{eq:psmr} 是觀測的對數死亡比 $\log(d_x/e_x)$ 與整體
$\log\mathrm{SMR}$ 的加權平均，其權重的設計使死亡數愈少的年齡愈倚賴參考母體。
此一形式與保險的信賴度估計同型 \citep{klugman2012}，
也與 \citet{kimeldorf1967} 的貝氏修勻同屬一類。

式~\eqref{eq:psmr} 有兩項性質值得指出。
首先，$d_x=0$ 時觀測端的權重為零，修勻值退化為 $\mathrm{SMR}\cdot m^R_x$，
恆為正。參考母體修勻因此根本不需要替代值，
零格的困難在這條路上是以外部資訊解決的，而非以慣例解決的。
該文亦即據此主張其方法「漸近不偏且不存在零計數問題」。
其次，$\hat h^2=0$ 時所有年齡都完全採用 $\mathrm{SMR}\cdot m^R_x$，
亦即整條年齡曲線的\textbf{形狀}來自參考母體，只有水準由目標區域決定。
這說明該法的表現必然取決於兩個母體的形狀是否相近，
而這一點與該文自己「標準母體的選擇不關鍵」的主張存在張力。

\citet{yue2019} 正是以此為主題。該文另考慮 \citet{whittaker1923} 修勻的
比值版本，即先取 $s_x=(d_x/P_x)/m^R_x$，再以
$\min_r \sum_x w_x(r_x-s_x)^2+h\sum_x(\Delta^z r_x)^2$ 求修勻後的比值，
取 $w_x$ 為曝露數、$h$ 為各年齡曝露數的平均。
該文設計七種比值情境（$s_x$ 固定為 $0.8$、$1.0$、$1.2$，以及隨年齡遞增、
遞減、V 型、倒 V 型），以 100,000 人對 2,000,000 人的模擬比較，
結論是：在前三種情境下 partial SMR 的平均絕對百分比誤差明顯最小
（約 $12\%$，不修勻為 $25$--$29\%$），但在遞增情境下惡化至 $47.65\%$；
Whittaker 比值修勻則在七種情境下都穩定在 $14$--$19\%$。
該文並指出，若把遞增情境的比值範圍由 $0.5$--$1.5$ 縮為 $(1-a)$--$(1+a)$，
則 $a\le 0.4$ 時 partial SMR 仍較佳。
綜上所述，這條脈絡的發展是：提出者以零格無法取對數為動機、以外部母體為解法，
並主張對該母體的選擇不敏感；後續的檢驗則界定了該主張成立的範圍。

合併母體與修勻兩條支線近年已被接在一起。\citet{pitt2018} 在廣義線性模型的
框架內同時施加 P-樣條修勻與線性約束，並以共同因子聯合配適兩性的死亡率，
亦即把「換成 Poisson」、「加修勻」與「借用另一個母體」三項補救放進同一個
目標函數求解，其結果在樣本外的預測誤差上優於逐項施加者。
這一條因而是第一類應對發展的終點：外部資訊能借的都借了。

這條脈絡所留下的缺口不在理論前提而在\textbf{資料的可得性}。
它要求分析者先有一個型態相近的參考母體，而在鄉鎮市區的應用中，
自然的參考母體是全國或鄰近地區，但臺灣縣市之間零歲平均餘命的最大差距
超過十歲 \citep{yue2019}，鄰近地區的社經組成亦常有明顯差異，
\citet{wangyue2026} 因而改以分群尋找合適的參考母體，
但同時指出相鄰地區的社經組成常有明顯差異，
使參考母體的選取本身成為一個待解的課題。
換言之，這條路把問題由「如何從少量資料估計」轉為「如何判斷參考母體是否
可用」，而後者無法由目標區域自己的資料回答。
下一小節所回顧的第二類應對即不需要外部母體，其代價則轉移到分配假設上。
